\section{Escalera de casos: células, organoides y retroalimentación de estructura y herramientas}

La clase continúa mostrando múltiples casos de colaboración. No deben reducirse a una frase como "IA descubrió muchos fármacos y proteínas", ya que cada caso tiene diferentes objetos de estudio, sistemas experimentales, estado de publicación y alcance extrapolable. Este capítulo lee AML, fibrosis hepática, ensamblaje inmune vegetal, diseño OCT4 y resultados tempranos de rejuvenecimiento por niveles de evidencia.

\subsection{Primero se construye un mapa de casos}

El diagrama general divide el trabajo en reutilización de fármacos contra el cáncer, reversión del daño hepático y descubrimiento e ingeniería de proteínas. El flujo común es que la IA propone u organiza candidatos y los equipos humanos experimentales seleccionan y validan; la diferencia radica en que las plataformas de validación varían desde líneas celulares y organoides hasta biología estructural, por lo que la fuerza de las conclusiones no puede compararse directamente de manera transversal.

\lecturefigure{slide-009.jpg}{Mapa de casos impulsados por colaboración humano--IA}{Grabación oficial Stanford CS25 V6 Lecture 07 00:42:00}

\begin{knowledgebox}{Lectura de la figura: tres direcciones comparten el flujo, no la intensidad de las conclusiones}
El lado izquierdo sobre AML se enfoca principalmente en cómo los fármacos afectan la viabilidad de líneas celulares cancerosas; el caso central de fibrosis hepática usa organoides hepáticos humanos para observar señales antifibroticas y regenerativas; el lado derecho vincula el cribado estructural con AlphaFold a experimentos de purificación y microscopía. Todos pueden probar que "los candidatos merecen continuar su investigación", pero están a distancias diferentes del tratamiento clínico, la eficacia a nivel de órgano y la confirmación funcional biológica.
\end{knowledgebox}

\begin{importantbox}{Tabla general de estado de los casos}
\begin{tabularx}{\textwidth}{L{0.20\textwidth}X X}
\toprule
Caso & Evidencia hasta la fecha en clase & Conclusión permitida por las apuntes \\
\midrule
AML / KIRA6 & Resultados in vitro en líneas celulares en arXiv v1 & Candidatos con citotoxicidad selectiva; requiere validación clínica adicional \\
Fibrosis hepática & bioRxiv v1, organoides hepáticos de origen humano & Dos modificadores epigenómicos muestran señales antifibroticas y regenerativas \\
Planta NRC7 & bioRxiv v1, cribado con AlphaFold, purificación, microscopía electrónica de negativo & SNI encontró ensamblaje atípico de 11-mero \\
OCT4 & Demostración de plausibilidad estructural en el apéndice de arXiv v1 & La estructura no se dañó significativamente; la función aún no está confirmada \\
Rejuvenecimiento & Diapositiva de clase, explícitamente sin publicar & Solo registrar resultados preliminares de clase; no sacar conclusiones definitivas \\
\bottomrule
\end{tabularx}
\end{importantbox}

\subsection{AML: Señales de KIRA6 y dos candidatos fallidos}

La tarea de reutilización de fármacos para la leucemia mieloide aguda (leucemia mieloide aguda, AML) requiere que el sistema busque nuevas indicaciones en fármacos existentes. El enfoque de la clase es KIRA6: un inhibidor IRE1$\alpha$ que muestra disminución dependiente de la dosis de viabilidad en varias líneas celulares de AML. La figura también conserva información clave negativa: dos otros "nuevos" fármacos propuestos por la IA no funcionaron.

\lecturefigure{slide-011.jpg}{Resultados de respuesta a dosis de KIRA6 en líneas celulares de AML}{Grabación oficial Stanford CS25 V6 Lecture 07 00:45:28; arXiv:2502.18864v1}

\begin{knowledgebox}{Lectura de la figura: primero mira la concentración en el eje horizontal, luego la posición de descenso en las distintas líneas celulares}
El eje horizontal es concentración logarítmica y el vertical viabilidad relativa. Una curva que desciende antes indica que una concentración menor puede suprimir la actividad celular; sin embargo, la sensibilidad no es igual entre todas las líneas celulares. Lo verdaderamente valioso es colocar la línea celular de AML y un control no maligno en el mismo diseño para buscar la ventana terapéutica, en lugar de mostrar solo una curva de descenso. La revisión de novedad en la figura también reconoce que ya hay evidencia de dianas relacionadas; la novedad radica en la combinación específica de fármaco e indicación.
\end{knowledgebox}

La forma común de Hill de respuesta a dosis es
$$
V(c)=V_{\min}+\frac{V_{\max}-V_{\min}}
{1+(c/IC_{50})^{h}}.
$$
Donde:
\begin{itemize}
  \item $c$ es la concentración del fármaco;
  \item $V(c)$ es la actividad celular medida a esa concentración;
  \item $IC_{50}$ es la concentración necesaria para que la actividad baje a nivel medio;
  \item $h$ es la pendiente de Hill, que describe lo empinada que es la curva;
  \item La $IC_{50}$ in vitro no puede sustituir directamente la dosis humana, la seguridad ni la eficacia.
\end{itemize}

\teachervoice{00:45:20--00:47:12, el docente señala activamente que dos otras sugerencias de nuevos fármacos no funcionan. Los resultados negativos no son un defecto del discurso, sino datos necesarios para estimar la tasa de aciertos del sistema, detectar sesgo de selección y diseñar el ranking de la siguiente ronda.}

\begin{warningbox}{Matar células no es tratamiento clínico}
Las líneas celulares carecen de farmacocinética, entorno inmune, distribución tisular, toxicidad y evolución de resistencia presentes en el cuerpo humano. Los resultados de KIRA6 apoyan continuar con estudios de mecanismo, animales y seguridad, pero no respaldan describirlo como una nueva terapia para AML ya confirmada.
\end{warningbox}

\subsection{Fibrosis hepática: Vorinostat en organoides humanos}

El caso de fibrosis hepática actualiza la plataforma de validación a multi-línea de organoides hepáticos humanos. El co-científico de IA recomienda modificadores epigenómicos y el equipo experimental evalúa secuencialmente la actividad antifibrotica y la toxicidad. Vorinostat es un inhibidor HDAC anticanceroso ya aprobado; en este modelo reduce los cambios estructurales del cromatina inducidos por TGF$\beta$, acompañados de señales antifibroticas y regenerativas de células hepáticas.

\lecturefigure{slide-012.jpg}{Resultados antifibroticos de Vorinostat en organoides hepáticos humanos}{Grabación oficial Stanford CS25 V6 Lecture 07 00:47:48; bioRxiv:2025.04.29.651320v1}

\begin{knowledgebox}{Lectura de la figura: la línea punteada es la base, no la línea de cura}
El eje vertical es el cambio en la actividad de fibroblastos relativo al grupo sin tratar (fold change), con la línea verde punteada aproximadamente en 1. El inducir fibrosis eleva la actividad; si un candidato fármaco traa la distribución hacia abajo o por debajo de la base, indica que bajo esas condiciones del organoide se reduce la actividad pro-fibrotica. Los diagramas de caja muestran la distribución de muestras, no un único punto óptimo. El valor de Vorinostat radica en que un fármaco conocido obtiene nuevo mecanismo y candidato a nueva indicación, pero la figura no incluye desenlaces humanos.
\end{knowledgebox}

Si se normaliza con la base sin tratar, el efecto antifibrotico puede escribirse como
$$
E_{\text{anti-fibrotic}}
=1-\frac{A_{\text{drug}}-A_{\text{baseline}}}
{A_{\text{fibrosis}}-A_{\text{baseline}}}.
$$
Donde:
\begin{itemize}
  \item $A_{\text{drug}}$ es el indicador de actividad fibrotica tras tratamiento con el candidato fármaco;
  \item $A_{\text{fibrosis}}$ es el indicador inducido a fibrosis sin administrar el candidato fármaco;
  \item $A_{\text{baseline}}$ es la base no inducida;
  \item Esta normalización ayuda a comparar efectos dentro del experimento, pero no elimina el error de extrapolación entre el modelo y el cuerpo humano.
\end{itemize}

\begin{knowledgebox}{Concepto previo: organoide}
Los organoides son modelos experimentales formados por auto-organización celular en un entorno tridimensional que replica parcialmente la estructura y función tisular. Se acercan más a las interacciones a nivel de tejido que las líneas celulares monocapa, pero generalmente aún carecen de flujo sanguíneo completo, inmunidad, nervios y metabolismo sistémico. Por tanto, la evidencia de organoides es más fuerte que el razonamiento basado solo en texto, pero sigue antecediendo a los estudios animales, preclínicos y humanos.
\end{knowledgebox}

\teachervoice{00:47:20--00:48:46, el docente considera la plataforma experimental como una comprobación de realidad biológica para las hipótesis de IA. La clase no empaqueta la mejora de organoides como un éxito de tratamiento en pacientes.}

\subsection{Estructura vegetal: cómo SNI convierte anomalías en candidatos}

Después de que AlphaFold genera miles de predicciones estructurales, los investigadores enfrentan otro problema: ¿cuáles resultados "que no parecen estructuras conocidas" son realmente novedosos y cuáles son simplemente errores de predicción? El equipo y el co-científico de IA diseñaron el Índice de Novedad Estructural (SNI), realizaron cribado estructural del ensamblaje de proteínas inmunes vegetales y entregaron candidatos anómalos a purificación y microscopía electrónica de negativo.

\lecturefigure{slide-013.jpg}{El Structural Novelty Index encuentra ensamblaje atípico de 11-mero en NRC7}{Grabación oficial Stanford CS25 V6 Lecture 07 00:50:32; bioRxiv:2026.05.03.722499v1}

\begin{knowledgebox}{Lectura de la figura: la IA reduce el espacio de búsqueda, los experimentos humanos deciden si la anomalía es real}
El mapa de calor superior compara los juicios de novedad estructural del científico y del co-científico de IA sobre los candidatos; el flujo inferior va desde el cribado SNI hacia expresión, purificación y microscopía electrónica, observándose finalmente un ensamblaje cíclico de NRC7 formado por 11 subunidades. El punto clave no es "que la IA ve mejor que el humano", sino que propone criterios adicionales para ayudar a seleccionar prioritariamente las anomalías dignas de costoso control experimental entre cientos de estructuras AlphaFold.
\end{knowledgebox}

\begin{importantbox}{Valor científico de detección de anomalías}
Cuando los modelos generadores producen candidatos masivos, lo más escaso no son los candidatos, sino la atención experimental. Indicadores como SNI convierten "la desviación de la arquitectura canónica" en una señal ordenable; luego se debe usar experimento independiente para distinguir novedad biológica del artefacto del modelo. El rol más valioso de un agente suele ser reducir el espacio de validación, no declarar que la anomalía ya está confirmada.
\end{importantbox}

\subsection{OCT4: poner AlphaFold en el bucle de retroalimentación}

OCT4 es uno de los factores Yamanaka. El co-científico de IA sugiere modificar la secuencia cerca del dominio POU, luego invoca a AlphaFold para verificar la plausibilidad estructural y devuelve los resultados estructurales al diseño de la siguiente ronda. El apéndice de arXiv v1 llama explícitamente a esto una demostración ilustrativa: para probar que las modificaciones mejoran la unión al ADN o mantienen la interacción con SOX2, se requieren afinidad de unión, efectos fuera de objetivo y validación en laboratorio húmedo.

\lecturefigure{slide-014.jpg}{Bucle de retroalimentación de diseño OCT4 compuesto por co-científico de IA y AlphaFold}{Grabación oficial Stanford CS25 V6 Lecture 07 00:51:16; arXiv:2502.18864v1 Apéndice A.6}

\begin{knowledgebox}{Lectura de la figura: pLDDT, pTM e ipTM solo responden a preguntas de confianza estructural}
La figura compara las estructuras predichas y los indicadores de confianza del OCT4 original y con secuencia modificada. pLDDT similar indica que la confianza en la estructura local no colapsó significativamente; pTM e ipTM ofrecen pistas sobre confiabilidad global del plegamiento y de la interfaz, respectivamente. No pueden probar directamente mejora de función transcripcional, cambio de destino celular o ausencia de efectos fuera de objetivo. La conclusión correcta es "valorable para continuar validación estructural", no "la proteína optimizada ya tiene éxito".
\end{knowledgebox}

\begin{lstlisting}[language=Python,caption={Flujo conceptual del proceso de retroalimentación de diseño de hipótesis donde entra el modelo especializado}]
sequence = co_scientist.propose_sequence(goal, constraints)
for iteration in range(max_rounds):
    structure = alphafold.predict(sequence)
    checks = assess(structure, metrics=["pLDDT", "pTM", "ipTM"])
    if checks.structurally_plausible:
        sequence = co_scientist.refine(sequence, checks)
    else:
        sequence = co_scientist.repair(sequence, checks.failures)
return send_to_binding_and_wet_lab_validation(sequence)
\end{lstlisting}

\begin{warningbox}{Uso de herramientas no es verdad de herramientas}
Invocar AlphaFold puede reducir secuencias manifiestamente irracionales, pero no convierte un modelo predictivo en instrumento experimental. Si el agente y la herramienta comparten sesgos, el bucle de retroalimentación puede optimizar erróneamente métricas proxy. Se debe añadir modelos heterogéneos, chequeo de base de datos, experimentos y condiciones de parada fuera del loop.
\end{warningbox}

\subsection{Rejuvenecimiento temprano: los resultados de clase deben mantener la línea roja}

La clase muestra resultados tempranos con descenso marcado en marcadores de senescencia en células de donantes ancianos, y el título mismo indica explícitamente `Sin publicar y requiere revisión por pares`. Estas figuras pueden crear fácilmente la impresión de que "la reversión del envejecimiento ya se logró"; por tanto, las apuntes deben registrar simultáneamente el objeto experimental, los marcadores, los controles y el hecho de que los métodos y la réplica aún no están publicados.

\lecturefigure{slide-015.jpg}{Resultados tempranos de clase para un factor de rejuvenecimiento nuevo}{Grabación oficial Stanford CS25 V6 Lecture 07 00:52:00; la diapositiva marca explícitamente sin publicar y requiere revisión por pares}

\begin{knowledgebox}{Lectura de la figura: la altura de las barras describe el marcador, no el flujo inverso de edad}
El eje vertical es la proporción de células juzgadas como senescentes; el grupo sin tratar de donantes ancianos es más alto, mientras que los grupos KLOTHO y dos factores descubiertos por IA muestran reducción marcada. La figura apoya una observación preliminar de "descenso del marcador bajo esas condiciones experimentales"; no indica que todas las funciones celulares se recuperaron a un estado joven, ni proporciona tamaño muestral, estadística completa, efectos fuera de objetivo o réplicas independientes.
\end{knowledgebox}

\begin{warningbox}{Sin publicar significa qué}
Sin publicar no equivale a error, ni a confirmación. Significa que el lector no puede verificar completamente los métodos, muestras, criterios de exclusión, análisis estadístico y réplica. Las apuntes conservan este caso para aprender sobre jerarquía de evidencia y dirección de investigación, sin escribirlo como una conclusión de terapia anti-envejecimiento.
\end{warningbox}

\teachervoice{00:51:40--00:53:32, el docente marca explícitamente "requiere revisión por pares" en la diapositiva. El tono del profesor es "resultados tempranos", no un lanzamiento oficial; cualquier lenguaje más fuerte cruzaría la evidencia de clase.}

\subsection{Resumen de la sección}

Los cinco casos forman una escalera desde la generación de candidatos hasta distintos niveles experimentales. AML proporciona respuesta a dosis en líneas celulares y expone candidatos fallidos; fibrosis hepática avanza a organoides humanos; planta NRC7 usa cribado estructural más purificación y confirmación por microscopía electrónica de 11-mero; OCT4 es solo retroalimentación de plausibilidad estructural; rejuvenecimiento sigue siendo un resultado de clase sin publicar. Un agente puede mejorar el descubrimiento de candidatos y la prioridad experimental, pero cada modelo experimental tiene límites claros de extrapolación.


```
